
\subsection{Land-Use Area Trends by Land Protection Pressure}

Descriptive analysis reveals a central contrast between high- and low-pressure prefectures. Following the 2017 tightening of cropland protection, urban expansion trajectories remain similar across prefectures, whereas cropland adjustment differs sharply by land protection pressure. Figure~\ref{fig:size} presents prefecture-demeaned trends in impervious surface area and cropland area, separately for core and non-core areas and by land protection pressure.

Across both core and non-core areas, trends in impervious surface expansion are broadly comparable between high- and low-pressure prefectures, with steady growth and no visible divergence around 2017. This pattern indicates that urban development proceeded along similar paths regardless of land protection pressure, in line with the continued prioritization of economic growth objectives.

In contrast, cropland dynamics diverge markedly by land protection pressure after 2017. Prior to the policy tightening, cropland area declines in both high- and low-pressure prefectures. After 2017, however, cropland area rebounds and increases in high-pressure prefectures, while continuing to decline and exhibiting only a limited and delayed recovery in low-pressure prefectures. This divergence is concentrated in non-core areas; cropland trends within core urban areas remain similar across prefecture types throughout the period.

These patterns suggest that the 2017 policy tightening primarily constrained high-pressure prefectures, while having limited effects in low-pressure prefectures. At the same time, urban expansion in core areas appears insulated from these constraints, implying that the burden of meeting binding cropland protection targets was shifted toward non-core areas within high-pressure prefectures.



\begin{figure}[h]
    \centering
    \includegraphics[width=0.8\linewidth]{Motivating_size.png}
    \caption{Trends in impervious area and cropland area by land type}
    \label{fig:size}
\end{figure}

\subsection{Vegetation Coverage over the Reclamation Cycle}

Figure~\ref{fig:fvc} illustrates the evolution of fractional vegetation coverage on newly reclaimed land as a function of years since reclamation. Pooling all newly reclaimed grids in the sample, mean FVC increases monotonically over time. This pattern indicates that vegetation coverage improves gradually following reclamation, indicating that FVC captures the intensity of agricultural use and land cultivation in the years after cropland creation.

Figure~\ref{fig:fvc_group} further disaggregates these patterns by land protection pressure and policy period. For grids reclaimed prior to 2018, as well as for grids located in low-pressure prefectures, FVC exhibits a similar upward trajectory to that observed in the full sample, increasing with years since reclamation. In contrast, newly reclaimed land in high-pressure prefectures after the policy tightening displays a markedly different pattern. These grids exhibit substantially lower initial FVC and show little subsequent improvement over time, with mean vegetation coverage remaining persistently low.

In a nutshell, these descriptive patterns suggest that while newly reclaimed land typically becomes increasingly cultivated over time, this process is disrupted for cropland created after the policy tightening in high-pressure prefectures. The persistently low vegetation coverage observed for these grids is consistent with a higher likelihood of underutilization or temporary fallow, whereas similar patterns are not observed in low-pressure prefectures. These motivating facts point to heterogeneous responses to the tightening of cropland protection, concentrated in regions where land protection constraints are most binding.


    \begin{figure}[htbp]
    \centering
    \includegraphics[width=0.8\linewidth]{Motivating_fvc.png}
    \caption{Evolution of Vegetation Coverage after Reclamation}
    \label{fig:fvc}
    \end{figure}

    \begin{figure}[htbp]
    \centering
    \includegraphics[width=0.8\linewidth]{Motivating_reclaimed_group.png}
    \caption{Vegetation Coverage after Reclamation by Land Protection Pressure and Polity Period}
    \label{fig:fvc_group}
    \end{figure}




